% Lösungen Einheit 4 von 5 – Umkehraufgaben (zusammenbau v0.1)
\einheitenkopf{Einheit 4 von 5 · Umkehraufgaben}

\erg{19}{a) n über den Logarithmus – Mindestanzahl bei „mindestens einmal“ \quad b) Nietenwahrscheinlichkeit $0{,}85$: $1 - 0{,}85^n \ge 0{,}9$ \quad c) $n = 11$: $1 - 0{,}64^n \ge 0{,}99 \Leftrightarrow 0{,}64^n \le 0{,}01 \Leftrightarrow n \ge \frac{\mathrm{ln}\, 0{,}01}{\mathrm{ln}\, 0{,}64} \approx 10{,}32$, aufrunden \quad d) höchstens $0{,}0522$ (etwa 5,2\,\%): $(1 - p)^{30} \ge 0{,}2 \Leftrightarrow 1 - p \ge \sqrt[30]{0{,}2}$ \quad e) $k = 27$: $P(X \le 26) \approx 0{,}9300 < 0{,}95$, $P(X \le 27) \approx 0{,}9588$ \quad f) $r = 8$: Erwartungswert 40; $P(32 \le X \le 48) \approx 0{,}9179$, bei r = 7 nur etwa 0{,}8747 \quad g) $n = 171$: $P(X > 120) \approx 0{,}8912$ für 170 und $\approx 0{,}9127$ für 171 (X: Anzahl mit Fahrrad) \quad h) $7\,\%$: $P(X \le 3) \approx 0{,}5327$ bei 7\,\%, $\approx 0{,}4253$ bei 8\,\% \quad i) Richtig: die Wahrscheinlichkeit für keinen Gewinn ist eine Potenz der Nietenwahrscheinlichkeit; mit jedem weiteren Los kommt ein Faktor kleiner als eins hinzu, sie fällt, also steigt die Wahrscheinlichkeit für mindestens einen Gewinn. \quad j) $p = \frac{1}{3}$, $n = 18$: $\frac{n p (1 - p)^{n-1}}{(1 - p)^n} = \frac{n p}{1 - p} = 9$, mit $n p = 6$ folgt $1 - p = \frac{2}{3}$}
\erg{20}{a) Richtig, es sind 5 Werte: i = 11 bis 15 (etwa 0{,}1021 bis 0{,}1110); die Nachbarn liegen darunter, $P(Y = 10) \approx 0{,}0749$ und $P(Y = 16) \approx 0{,}0867$}
\erg{21}{a) Fehler: $\mathrm{ln}\, 0{,}85$ ist negativ, beim Teilen dreht sich das Zeichen: $n \ge 14{,}17$. Richtig: mindestens $n = 15$ Lose \quad b) Sie ist eine Potenz der Nietenwahrscheinlichkeit; diese liegt zwischen null und eins, und mit jedem Versuch wird mit ihr multipliziert – das Produkt wird kleiner. \quad c) $3$ Melder: $0{,}1^n \le 0{,}001$ gilt ab $n = 3$}
